Logs and data exports carry personal data and credentials and often exceed memory. Chunked processing adds a silent failure: a value straddling a chunk boundary is split, missed and written out in the clear. SafeStream-Redactor treats chunking as a correctness problem. Each detector declares its longest span and the context it reads. We show that an overlap of the longest span can still leak, and prove that if the overlap covers every detector's span plus its right context, the streamed output equals a whole-document pass, parallel output equals sequential output or falls back to it, and the working buffer is independent of input size; the engine enforces this overlap. Naive 1\,KiB chunking leaked 237 RedactionBench values a whole-document pass redacts and missed 388 in real logs; overlapping windows still missed some, SafeStream none. Streaming added 8.4--8.54\,MiB from 8\,MiB to 2\,GiB, against 904.6\,MiB for whole-document redaction of 32\,MiB. On 14,626 real credentials labelled by people, the pattern tiers fully redacted 0.55 against 0.40 for gitleaks, robust to resampling by repository, at a higher false-alarm rate. An opt-in tier caught all 618 user names in held-out Linux logs. On 1,014 held-out court judgments, a free-text configuration frozen before scoring reached character F1 0.67 against 0.54 for Presidio (0.76 against 0.63 on Wikipedia biographies) and fully redacted 89\% of direct identifiers against 13\%. Yet no configuration fully protected more than 17.9\% of RedactionBench documents.
